\chapter{Introduction générale}

\section{Contexte général}
Le secteur du développement logiciel et des applications mobiles connaît 
aujourd'hui une croissance exponentielle, portée par la digitalisation 
massive des services et l'adoption généralisée des smartphones dans notre 
quotidien. Dans ce contexte hautement concurrentiel, la qualité de 
l'expérience utilisateur est devenue un facteur différenciant majeur 
pour toute entreprise proposant une application mobile. Les utilisateurs 
sont de plus en plus exigeants et n'hésitent pas à abandonner une 
application qui ne répond pas à leurs attentes ou qui ne leur offre 
pas un canal de communication direct avec l'équipe produit.

C'est dans ce contexte que s'inscrit mon stage de licence au sein de 
Bewize, une entreprise innovante proposant une application mobile à 
destination de ses utilisateurs. Bewize, consciente de l'importance 
de l'écoute client dans l'amélioration continue de son produit, a 
identifié un besoin crucial : permettre à ses utilisateurs de transmettre 
facilement leurs feedbacks et réclamations directement depuis l'application 
mobile, et centraliser ces retours dans un espace dédié accessible 
par l'équipe interne.

\section{Présentation de l'entreprise Bewize}
Bewize est une entreprise spécialisée dans le développement d'applications 
mobiles, dont le siège est situé à Casablanca. Elle développe et maintient 
une application mobile destinée à ses utilisateurs, accompagnée d'un 
dashboard monitor interne permettant à ses équipes de gérer et analyser 
les données liées à l'utilisation de l'application.

L'équipe de Bewize est organisée autour de plusieurs pôles complémentaires : 
une équipe produit pilotée par le Product Manager, une équipe technique 
encadrée par le Tech Lead, une équipe de développement et une équipe 
marketing digital. Cette organisation pluridisciplinaire favorise une 
approche holistique du développement produit, où les considérations 
techniques, fonctionnelles et marketing sont constamment alignées.

Mon stage s'est déroulé au sein de cette équipe dynamique et bienveillante, 
sous la supervision directe du Product Manager et avec l'encadrement 
technique du Tech Lead. Cette double supervision m'a permis de bénéficier 
à la fois d'une vision produit claire et d'un accompagnement technique 
de qualité tout au long de mon stage.

\section{Problématique}
Avant le développement de la fonctionnalité objet de ce stage, les 
utilisateurs de l'application mobile Bewize ne disposaient d'aucun 
moyen simple et intégré pour communiquer leurs retours à l'équipe Bewize. 
Cette absence de canal de communication directe engendrait plusieurs 
problèmes concrets :

\begin{itemize}
    \item Les utilisateurs insatisfaits n'avaient pas de moyen simple 
    pour signaler un problème ou exprimer une réclamation, ce qui 
    pouvait conduire à l'abandon de l'application sans que l'équipe 
    Bewize ne soit informée des causes.
    
    \item L'équipe produit ne disposait pas d'une source structurée 
    et centralisée de retours utilisateurs pour orienter ses décisions 
    de développement et améliorer continuellement le produit.
    
    \item Le traitement des rares réclamations reçues par d'autres 
    canaux (email, réseaux sociaux) était manuel, non systématisé 
    et difficile à suivre dans le temps.
    
    \item L'absence de données structurées sur la satisfaction 
    utilisateur rendait difficile l'évaluation objective de la qualité 
    de l'application et l'identification des axes d'amélioration prioritaires.
\end{itemize}

Face à cette problématique, la question centrale de mon stage a été 
la suivante : \textbf{Comment développer une fonctionnalité simple, 
accessible et fiable permettant aux utilisateurs de l'application 
mobile Bewize de transmettre facilement leurs feedbacks et réclamations, 
et à l'équipe interne de les consulter et les traiter efficacement ?}

\section{Objectifs du stage}
Pour répondre à cette problématique, mon stage chez Bewize s'est 
articulé autour de quatre objectifs principaux, définis conjointement 
avec le Product Manager lors de la phase de cadrage :

\begin{itemize}
    \item \textbf{Objectif 1 — Développer le formulaire WebView :} 
    concevoir et développer une page WebView en React JS contenant 
    un formulaire simple et intuitif permettant à l'utilisateur de 
    soumettre un feedback ou une réclamation en quelques clics depuis 
    l'application mobile.
    
    \item \textbf{Objectif 2 — Intégrer dans l'application mobile :} 
    intégrer la page WebView dans l'application mobile React Native 
    existante de Bewize, en ajoutant un bouton d'accès dédié et en 
    assurant une navigation fluide entre les différents écrans.
    
    \item \textbf{Objectif 3 — Développer le back-end :} créer les 
    APIs REST nécessaires avec Spring Boot pour recevoir, valider, 
    traiter et persister les données de feedbacks et réclamations 
    dans la base de données PostgreSQL.
    
    \item \textbf{Objectif 4 — Développer le dashboard monitor :} 
    enrichir le dashboard monitor interne de Bewize avec une interface 
    de consultation et de traitement des feedbacks et réclamations 
    reçus, avec des fonctionnalités de filtrage et de gestion des statuts.
\end{itemize}

\section{Méthodologie adoptée}
Pour atteindre ces objectifs dans les délais impartis du stage, 
j'ai travaillé en suivant la méthodologie Agile Scrum, déjà adoptée 
par l'équipe de développement de Bewize. Cette approche itérative 
et incrémentale m'a permis d'organiser mon travail en trois Sprints 
de deux semaines chacun, de livrer régulièrement des résultats 
tangibles au Product Manager et de m'adapter rapidement aux 
ajustements demandés en cours de développement.

Les outils et technologies utilisés pour mener à bien ce projet 
sont les suivants : React JS et React Native pour le développement 
front-end, Spring Boot pour le back-end, PostgreSQL pour la base 
de données, Docker pour la conteneurisation, et Git avec GitHub 
pour le versioning et la collaboration.

\section{Structure du rapport}
Ce rapport de stage est organisé en cinq chapitres qui suivent 
la progression naturelle du projet, depuis l'analyse des besoins 
jusqu'aux résultats obtenus :

\begin{itemize}
    \item \textbf{Chapitre 1 — Présentation de l'entreprise :} 
    ce chapitre présente Bewize, son activité, son organisation 
    et son écosystème technique existant avant le début du stage.
    
    \item \textbf{Chapitre 2 — Cadrage de besoin et analyse :} 
    ce chapitre détaille l'analyse des besoins fonctionnels et 
    non fonctionnels, l'étude de l'existant, la solution proposée 
    et le planning du stage.
    
    \item \textbf{Chapitre 3 — Conception et architecture :} 
    ce chapitre présente la méthodologie Agile Scrum adoptée, 
    les technologies utilisées, l'architecture globale de la 
    solution et la modélisation UML du système.
    
    \item \textbf{Chapitre 4 — Mise en œuvre :} ce chapitre 
    décrit en détail l'implémentation technique de chaque composant 
    de la solution, depuis le formulaire WebView jusqu'au dashboard 
    monitor, en passant par le back-end Spring Boot.
    
    \item \textbf{Chapitre 5 — Test et résultats :} ce chapitre 
    présente la stratégie de test adoptée, les tests réalisés 
    à différents niveaux et les résultats obtenus, confirmant 
    la qualité et la fiabilité de la solution développée.
\end{itemize}

Ce rapport se conclut par une conclusion générale qui dresse 
le bilan du projet, présente les compétences acquises et propose 
des perspectives d'évolution pour la solution développée.